\chapter{Semigruppi, monoidi e gruppi}\label{cap:gruppi}

Una \textbf{struttura algebrica} è un insieme dotato di una o più operazioni. In questo capitolo studiamo le
strutture con \textbf{una sola operazione binaria}: semigruppi, monoidi e gruppi, le loro sottostrutture e le funzioni
che ne rispettano l'operazione (gli omomorfismi).

% ──────────────────────────────────────────────────────────────
\section{Operazioni e strutture algebriche}\label{sec:operazioni-interne}

\dfn[operazione-interna]{Operazione interna}{
    Sia $s$ un insieme non vuoto. Una funzione
    \[ * \colon \underbrace{s \times \dots \times s}_{n \text{ volte}} \to s \]
    si dice \textbf{operazione $n$-aria interna} di $s$. Per $n = 2$ si parla di \textbf{operazione binaria interna}
    e, se $x, y \in s$, scriveremo $x * y := *(x, y)$, abbandonando la notazione funzionale.
}

\dfn[commutativa-associativa]{Operazione commutativa e associativa}{
    Un'operazione binaria interna $*$ di $s$ si dice
    \begin{enumerate}
        \item \textbf{commutativa} se $\forall a, b \in s\,(a * b = b * a)$;
        \item \textbf{associativa} se $\forall a, b, c \in s\,\big(a * (b * c) = (a * b) * c\big)$.
    \end{enumerate}
}

\dfn[struttura-algebrica]{Struttura algebrica}{
    Siano $*_1, \dots, *_n$ operazioni interne di $s$. La $(n+1)$-upla $(s, *_1, \dots, *_n)$ si dice
    \textbf{struttura algebrica} ad $n$ operazioni. In particolare $(s, *)$, con $*$ binaria, è una struttura algebrica
    ad un'operazione binaria interna.
}

\dfn[semigruppo]{Semigruppo}{
    Sia $(s, *)$ una struttura algebrica con un'operazione binaria interna. $(s, *)$ si dice \textbf{semigruppo} se $*$ è
    associativa; si dice \textbf{semigruppo commutativo} se $*$ è sia associativa che commutativa.
}

% ──────────────────────────────────────────────────────────────
\section{Elemento neutro e monoidi}\label{sec:neutro}

\dfn[neutro]{Elemento neutro}{
    Sia $(s, *)$ una struttura algebrica con un'operazione binaria interna.
    \begin{itemize}
        \item $l \in s$ tale che $\forall x \in s\,(l * x = x)$ si dice \textbf{elemento neutro sinistro};
        \item $r \in s$ tale che $\forall x \in s\,(x * r = x)$ si dice \textbf{elemento neutro destro};
        \item $u \in s$ che sia neutro sia a sinistra sia a destra si dice \textbf{elemento neutro}.
    \end{itemize}
}

\begin{concetto}
\begin{Teorema}{Unicità dell'elemento neutro}{unicita-neutro}
    \begin{prerequisiti}
        \dip{Definizione}{def:neutro}
        \dip{Definizione}{def:semigruppo}
    \end{prerequisiti}
    Sia $(s, *)$ un semigruppo. Siano $l$ neutro a sinistra e $r$ neutro a destra. Allora $l = r$.
    In particolare, il neutro, se esiste, è unico.
\end{Teorema}
\begin{dimostrazione}
    \begin{caso}{$l = r$}
        $l = l * r = r$: la prima uguaglianza perché $r$ è neutro a destra, la seconda perché $l$ è neutro a sinistra.
    \end{caso}
    \begin{caso}{Il neutro è unico}
        Se $u$ e $u'$ sono due elementi neutri, $u$ è in particolare neutro a sinistra e $u'$ neutro a destra, quindi $u = u'$.
    \end{caso}
\end{dimostrazione}
\end{concetto}

\dfn[monoide]{Monoide}{
    Un semigruppo $(m, *)$ dotato di elemento neutro si dice \textbf{monoide}; per evidenziare il neutro scriveremo
    $(m, *, 1_m)$, dove $1_m$ è l'elemento neutro. Se $*$ è anche commutativa, il monoide si dice \textbf{commutativo}.
}

\begin{esempio}{Monoidi e non}
    \begin{itemize}
        \item $(\mbN, +, 0)$, $(\mbN, \cdot, 1)$, $(\mbZ, +, 0)$, $(\mbZ, \cdot, 1)$ sono tutti monoidi commutativi.
        \item Dato un insieme $a$, $(\mcP(a), \cap, a)$, $(\mcP(a), \cup, \varnothing)$ e $(\mcP(a), \Delta, \varnothing)$ sono monoidi
              commutativi, per le proprietà dimostrate nella Proposizione~\ref{prop:prop-operazioni} (e perché per ogni $x \subseteq a$
              si ha $x \cap a = x$, $x \cup \varnothing = x$, $x \,\Delta\, \varnothing = x$).
        \item $(\mbZ, -)$ \textbf{non} è un semigruppo: $(1 - 1) - 1 = -1 \neq 1 = 1 - (1 - 1)$. Ha $0$ come elemento neutro destro
              ($x - 0 = x$), che però non è neutro sinistro ($0 - x = -x$).
        \item Per un qualunque insieme non vuoto $a$, $(\mathrm{In}(a, a), \circ, \mathrm{id}_a)$, $(\mathrm{Su}(a, a), \circ, \mathrm{id}_a)$ e
              $(\mathrm{Bi}(a, a), \circ, \mathrm{id}_a)$ sono monoidi (Teoremi~\ref{teo:comp-iniettive}, \ref{teo:comp-suriettive},
              Corollario~\ref{cor:comp-biettive} e Proposizione~\ref{prop:assoc-prodotto}).
    \end{itemize}
\end{esempio}

\dfn[parole]{Monoide delle parole}{
    Sia $A$ un insieme di simboli (ad esempio $A = \{a, b, c\}$), detto \textbf{alfabeto}. L'insieme $s$ delle \textbf{parole} su $A$
    è l'insieme delle successioni finite di elementi di $A$, compresa la parola vuota $\varnothing$:
    \[ s = \{x \mid \exists n \in \mbN \setminus \{0\}\ \exists x_1, \dots, x_n \in A\,(x = x_1 \dots x_n)\} \cup \{\varnothing\}. \]
    La \textbf{concatenazione} $* \colon (x, y) \in s \times s \mapsto xy \in s$ (la parola $x$ seguita dalla parola $y$), con la
    convenzione $\varnothing * x = x = x * \varnothing$, rende $(s, *, \varnothing)$ un monoide.
}

\begin{concetto}
\begin{Proposizione}{Quando il monoide delle parole è commutativo}{parole-commutativo}
    \begin{prerequisiti}
        \dip{Definizione}{def:parole}
        \dip{Definizione}{def:commutativa-associativa}
    \end{prerequisiti}
    Sia $A$ un alfabeto non vuoto. Il monoide delle parole su $A$ è commutativo se e solo se $A = \{a\}$ ha un solo elemento.
\end{Proposizione}
\begin{dimostrazione}
    \begin{caso}{$(\leftarrow)$ Se $A = \{a\}$, allora il monoide è commutativo}
        Le parole sono formate solo dalla lettera $a$: $x = a \dots a$ ($n$ volte) e $y = a \dots a$ ($k$ volte), quindi sia
        $xy$ sia $yx$ sono la parola formata da $n + k$ lettere $a$.
    \end{caso}
    \begin{caso}{$(\rightarrow)$ Se il monoide è commutativo, allora $A$ ha un solo elemento}
        Per contrapposizione: se $A$ contiene almeno due lettere distinte $a$ e $b$, allora $a * b = ab \neq ba = b * a$,
        quindi il monoide non è commutativo.
    \end{caso}
\end{dimostrazione}
\end{concetto}

\dfn[duale]{Operazione duale (o opposta)}{
    Sia $(s, *)$ una struttura algebrica con un'operazione binaria interna. Definisco
    \[ \bar{*} \colon (a, b) \in s \times s \mapsto b * a \in s; \]
    $\bar{*}$ si dice \textbf{operazione duale} (o \textbf{opposta}) di $*$.
}

\begin{concetto}
\begin{Proposizione}{Proprietà dell'operazione duale}{duale}
    \begin{prerequisiti}
        \dip{Definizione}{def:duale}
        \dip{Definizione}{def:commutativa-associativa}
        \dip{Definizione}{def:neutro}
        \dip{Teorema}{teo:biettiva-inversa}
    \end{prerequisiti}
    Sia $(s, *)$ una struttura algebrica con un'operazione binaria interna.
    \begin{enumerate}
        \item $*$ è commutativa se e solo se lo è $\bar{*}$; $*$ è associativa se e solo se lo è $\bar{*}$.
        \item $l$ è neutro a sinistra in $(s, *)$ se e solo se è neutro a destra in $(s, \bar{*})$ (e viceversa).
        \item $\bar{\bar{*}} = *$; quindi la funzione $* \mapsto \bar{*}$, dall'insieme $\mathrm{Map}(s \times s, s)$ delle operazioni
              binarie interne di $s$ in se stesso, è l'inversa di sé stessa ed è biettiva.
    \end{enumerate}
\end{Proposizione}
\begin{dimostrazione}
    \begin{caso}{Punto 1}
        Commutatività: $a \,\bar{*}\, b = b \,\bar{*}\, a \leftrightarrow b * a = a * b$. Associatività: $a \,\bar{*}\, (b \,\bar{*}\, c) = (c * b) * a$
        e $(a \,\bar{*}\, b) \,\bar{*}\, c = c * (b * a)$, quindi l'associatività di $\bar{*}$ per $a, b, c$ è l'associatività di $*$ per $c, b, a$.
    \end{caso}
    \begin{caso}{Punto 2}
        $\forall x\,(l * x = x) \leftrightarrow \forall x\,(x \,\bar{*}\, l = x)$, per definizione di $\bar{*}$.
    \end{caso}
    \begin{caso}{Punto 3}
        $a \,\bar{\bar{*}}\, b = b \,\bar{*}\, a = a * b$. Una funzione che è l'inversa di sé stessa è biettiva (Teorema~\ref{teo:biettiva-inversa}).
    \end{caso}
\end{dimostrazione}
\end{concetto}

% ──────────────────────────────────────────────────────────────
\section{Elementi invertibili e gruppi}\label{sec:invertibili}

\dfn[invertibile]{Elemento invertibile e inverso (simmetrico)}{
    Sia $(m, *, 1_m)$ un monoide e sia $x \in m$.
    \begin{enumerate}
        \item Se esiste $l \in m$ con $l * x = 1_m$, $x$ lo dico \textbf{invertibile a sinistra} ed $l$ è detto
              \textbf{inverso (simmetrico) sinistro} di $x$.
        \item Se esiste $r \in m$ con $x * r = 1_m$, $x$ lo dico \textbf{invertibile a destra} ed $r$ è detto
              \textbf{inverso (simmetrico) destro} di $x$.
    \end{enumerate}
    Un elemento che è inverso sia destro sia sinistro di $x$ si dice \textbf{inverso} (o \textbf{simmetrico}) di $x$, e $x$ si dice
    \textbf{invertibile} (o \textbf{simmetrizzabile}).
}

\begin{concetto}
\begin{Teorema}{Unicità dell'inverso}{unicita-inverso}
    \begin{prerequisiti}
        \dip{Definizione}{def:invertibile}
        \dip{Definizione}{def:monoide}
    \end{prerequisiti}
    Sia $(m, *, 1_m)$ un monoide e sia $x \in m$ con un inverso sinistro $l$ e un inverso destro $r$. Allora $l = r$.
    In particolare, l'inverso, se esiste, è unico.
\end{Teorema}
\begin{dimostrazione}
    \begin{caso}{$l = r$}
        Usando l'associatività:
        \[ l = l * 1_m = l * (x * r) = (l * x) * r = 1_m * r = r. \]
    \end{caso}
    \begin{caso}{L'inverso è unico}
        Se $y$ e $y'$ sono due inversi di $x$, $y$ è in particolare inverso sinistro e $y'$ inverso destro, quindi $y = y'$.
    \end{caso}
\end{dimostrazione}
\end{concetto}

\info{Notazione}{
    L'inverso di $x$ lo scrivo $x^{-1}$. Se l'operazione si indica con $+$ (notazione additiva, ad esempio in $(\mbZ, +, 0)$)
    si scrive $-x$ e si chiama \textbf{opposto} di $x$.
}

\dfn[gruppo]{Gruppo, gruppo abeliano}{
    Un monoide in cui tutti gli elementi sono invertibili si chiama \textbf{gruppo}. Un gruppo commutativo si chiama
    \textbf{gruppo abeliano}.
}

\dfn[unita]{Insieme degli elementi invertibili}{
    Sia $(m, *, 1_m)$ un monoide. Scrivo $U(m)$ l'insieme di tutti gli \textbf{elementi invertibili} di $m$.
}

\begin{concetto}
\begin{Proposizione}{Gli invertibili formano un gruppo}{unita-gruppo}
    \begin{prerequisiti}
        \dip{Definizione}{def:unita}
        \dip{Teorema}{teo:unicita-inverso}
        \dip{Definizione}{def:gruppo}
    \end{prerequisiti}
    Sia $(m, *, 1_m)$ un monoide e siano $a, b \in U(m)$. Allora $a * b \in U(m)$ e $(a * b)^{-1} = b^{-1} * a^{-1}$.
    Di conseguenza $U(m)$, con l'operazione $*$, è un gruppo.
\end{Proposizione}
\begin{dimostrazione}
    \begin{caso}{$b^{-1} * a^{-1}$ è l'inverso di $a * b$}
        $(a * b) * (b^{-1} * a^{-1}) = a * (b * b^{-1}) * a^{-1} = a * a^{-1} = 1_m$ e, allo stesso modo,
        $(b^{-1} * a^{-1}) * (a * b) = b^{-1} * (a^{-1} * a) * b = b^{-1} * b = 1_m$.
    \end{caso}
    \begin{caso}{$U(m)$ è un gruppo}
        L'operazione resta associativa, $1_m \in U(m)$ (perché $1_m * 1_m = 1_m$) e ogni $a \in U(m)$ ha inverso $a^{-1}$, che è a sua volta
        invertibile (con inverso $a$), quindi sta in $U(m)$.
    \end{caso}
\end{dimostrazione}
\end{concetto}

\begin{esempio}{Elementi invertibili}
    $U(\mbN, +, 0) = \{0\}$; \ $U(\mbZ, \cdot, 1) = \{-1, 1\}$; \ se $a$ è un insieme non vuoto, $U(\mathrm{In}(a, a), \circ, \mathrm{id}_a) = \mathrm{Bi}(a, a)$
    (Teorema~\ref{teo:biettiva-inversa}), quindi $(\mathrm{Bi}(a, a), \circ)$ è un gruppo.
\end{esempio}

% ──────────────────────────────────────────────────────────────
\section{Parti chiuse e sottostrutture}\label{sec:sottostrutture}

\dfn[parte-chiusa]{Parte chiusa}{
    Sia $(s, *)$ una struttura algebrica con un'operazione binaria interna e sia $t \subseteq s$. $t$ si dice
    \textbf{parte chiusa} (o \textbf{stabile}) di $(s, *)$ se
    \[ \forall x, y \in t\,(x * y \in t). \]
}

\dfn[sottostruttura]{Operazione indotta, sottostruttura}{
    Se $t$ è una parte chiusa di $(s, *)$, definisco
    \[ *' \colon (x, y) \in t \times t \mapsto x * y \in t, \]
    detta \textbf{operazione indotta da $*$ su $t$}. $(t, *')$ la dico \textbf{sottostruttura} di $(s, *)$ e scrivo $(t, *') \leq (s, *)$
    (spesso, con abuso di notazione, $(t, *)$ o semplicemente $t \leq s$). Inoltre:
    \begin{itemize}
        \item se $(s, *)$ è un semigruppo, una sua sottostruttura si dice \textbf{sottosemigruppo};
        \item se $(m, *, 1_m)$ è un monoide, un \textbf{sottomonoide} è una sottostruttura che contiene $1_m$;
        \item se $(g, *)$ è un gruppo, un \textbf{sottogruppo} è un sottomonoide che contiene l'inverso di ogni suo elemento.
    \end{itemize}
}

\begin{esempio}{Sottostrutture}
    \begin{itemize}
        \item $\{0\}$ è un sottosemigruppo e un sottomonoide di $(\mbN, +, 0)$.
        \item $\{1\}$ è un sottomonoide di $(\mbN, \cdot, 1)$, ma non è una parte chiusa di $(\mbN, +)$: $1 + 1 = 2 \notin \{1\}$.
        \item Per $m \in \mbN$, $m\mbN = \{n \in \mbN \mid \exists k \in \mbN\,(n = km)\}$ è un sottomonoide di $(\mbN, +, 0)$.
        \item Per $m \in \mbZ$, $m\mbZ = \{n \in \mbZ \mid \exists k \in \mbZ\,(n = km)\}$, l'insieme dei multipli interi di $m$, è un
              sottogruppo di $(\mbZ, +)$ e $(m\mbZ, \cdot)$ è un sottosemigruppo di $(\mbZ, \cdot)$. Ad esempio $0\mbZ = \{0\}$ e $1\mbZ = \mbZ$.
        \item Se $a$ è un insieme non vuoto, $(\mathrm{Bi}(a, a), \circ)$ è un sottomonoide di $(\mathrm{In}(a, a), \circ)$, ed è un gruppo.
    \end{itemize}
\end{esempio}

\begin{concetto}
\begin{Teorema}{Intersezione di parti chiuse}{intersezione-chiuse}
    \begin{prerequisiti}
        \dip{Definizione}{def:parte-chiusa}
        \dip{Definizione}{def:sottostruttura}
        \dip{Definizione}{def:intersezione}
    \end{prerequisiti}
    Sia $(s, *)$ una struttura e $f$ una famiglia di parti chiuse di $(s, *)$. Allora $\bigcap f$ è una parte chiusa di $(s, *)$.
    Inoltre, se $f$ è una famiglia (non vuota) di sottomonoidi (risp.\ sottogruppi), $\bigcap f$ è un sottomonoide (risp.\ sottogruppo).
\end{Teorema}
\begin{dimostrazione}
    \begin{caso}{$\bigcap f$ è una parte chiusa}
        Siano $x, y \in \bigcap f$; allora $\forall t \in f\,(x, y \in t)$. Ma per ipotesi ogni $t \in f$ è una parte chiusa, quindi
        $\forall t \in f\,(x * y \in t)$, cioè $x * y \in \bigcap f$.
    \end{caso}
    \begin{caso}{Sottomonoidi e sottogruppi}
        Se ogni $t \in f$ contiene $1_m$, anche $\bigcap f$ lo contiene. Se ogni $t \in f$ contiene l'inverso di ogni suo elemento e
        $x \in \bigcap f$, allora $x^{-1}$ sta in ogni $t \in f$, quindi $x^{-1} \in \bigcap f$.
    \end{caso}
\end{dimostrazione}
\end{concetto}

\dfn[generata]{Sottostruttura generata}{
    Sia $(s, *)$ una struttura algebrica e sia $t \subseteq s$. Dico \textbf{sottostruttura di $(s, *)$ generata da $t$}
    \[ \langle t \rangle := \bigcap \{c \in \mcP(s) \mid c \leq s \wedge t \subseteq c\}, \]
    l'intersezione delle sottostrutture di $s$ che contengono $t$ (dello stesso tipo: sottomonoidi se $s$ è un monoide, sottogruppi
    se $s$ è un gruppo). Per il Teorema~\ref{teo:intersezione-chiuse} è anch'essa una sottostruttura: è la \textbf{più piccola}
    sottostruttura di $s$ che contiene $t$. Si scrive $\langle x \rangle$ invece di $\langle \{x\} \rangle$, $\langle a, b \rangle$ invece di
    $\langle \{a, b\} \rangle$. Una struttura $s$ si dice \textbf{ciclica} se esiste $x \in s$ con $s = \langle x \rangle$.
}

\begin{esempio}{Sottostrutture generate}
    \begin{itemize}
        \item Se $(m, *, 1_m)$ è un monoide, $\langle \varnothing \rangle = \{1_m\}$.
        \item Il sottomonoide di $(\mbN, +, 0)$ generato da $1$ è $\mbN$; quello generato da $2$ è $2\mbN$.
        \item Il sottosemigruppo di $(\mbZ, +)$ generato da $2$ è $2\mbN \setminus \{0\}$.
        \item Il sottogruppo di $(\mbZ, +)$ generato da $2$ è $2\mbZ$; quello generato da $\{2, 3\}$ è $\mbZ$.
        \item $(\mbN, +, 0)$ è un monoide ciclico ($\mbN = \langle 1 \rangle$) e $(\mbZ, +)$ è un gruppo ciclico ($\mbZ = \langle 1 \rangle$).
    \end{itemize}
\end{esempio}

\begin{concetto}
\begin{Teorema}{Caratterizzazione dei sottomonoidi generati}{sottomonoidi-generati}
    \begin{prerequisiti}
        \dip{Definizione}{def:generata}
        \dip{Teorema}{teo:intersezione-chiuse}
        \dip{Proposizione}{prop:doppia-inclusione}
        \dip{Definizione}{def:sottostruttura}
    \end{prerequisiti}
    Sia $(m, *, 1_m)$ un monoide e sia $t \in \mcP(m) \setminus \{\varnothing\}$. Allora
    \[ \langle t \rangle = \{x \in m \mid \exists n \in \mbN \setminus \{0\}\ \exists x_1, \dots, x_n \in t\,(x = x_1 * \dots * x_n)\} \cup \{1_m\}, \]
    cioè l'insieme delle «combinazioni» (prodotti finiti) di elementi di $t$, più il neutro.
\end{Teorema}
\begin{dimostrazione}
    Chiamo $S$ l'insieme a destra e, tramite la doppia inclusione, dimostro che $\langle t \rangle = S$. Ricordiamo che
    $\langle t \rangle = \bigcap\{c \in \mcP(m) \mid t \subseteq c \wedge c \leq m\}$.
    \begin{caso}{$(\subseteq)$ $\langle t \rangle \subseteq S$}
        Basta vedere che $S$ è un sottomonoide che contiene $t$: allora è uno degli insiemi che si intersecano, e l'intersezione
        è contenuta in esso.
        Se prendo $n = 1$, ho $t \subseteq S$; inoltre $1_m \in S$. Prendo $x_1 * \dots * x_{n_1}$ e $y_1 * \dots * y_{n_2}$ in $S$: allora
        $x_1 * \dots * x_{n_1} * y_1 * \dots * y_{n_2} \in S$ (è ancora un prodotto finito di elementi di $t$; se uno dei due è $1_m$ il prodotto
        è l'altro). Quindi $S \leq m$ e $\langle t \rangle \subseteq S$.
    \end{caso}
    \begin{caso}{$(\supseteq)$ $S \subseteq \langle t \rangle$}
        Prendo un qualunque sottomonoide $e \leq m$ con $t \subseteq e$. Allora $1_m \in e$, e se $x_1, \dots, x_n \in t$ allora
        $x_1, \dots, x_n \in e$, quindi $x_1 * \dots * x_n \in e$ poiché $e$ è una parte chiusa. Quindi $S \subseteq e$; per la generalità
        di $e$, $S$ è contenuto nell'intersezione, cioè $S \subseteq \langle t \rangle$.
    \end{caso}
\end{dimostrazione}
\end{concetto}

\begin{concetto}
\begin{Teorema}{Caratterizzazione dei sottogruppi generati}{sottogruppi-generati}
    \begin{prerequisiti}
        \dip{Teorema}{teo:sottomonoidi-generati}
        \dip{Definizione}{def:gruppo}
    \end{prerequisiti}
    Sia $(g, *)$ un gruppo e sia $t \in \mcP(g) \setminus \{\varnothing\}$. Allora
    \[ \langle t \rangle = \{x \in g \mid \exists n \in \mbN \setminus \{0\}\ \exists \varepsilon_1, \dots, \varepsilon_n \in \{-1, 1\}\ \exists x_1, \dots, x_n \in t\,(x = x_1^{\varepsilon_1} * \dots * x_n^{\varepsilon_n})\} \cup \{1_g\}. \]
\end{Teorema}
\begin{dimostrazione}[Nota]
    Negli appunti questo teorema è dato senza dimostrazione. L'idea è la stessa del Teorema~\ref{teo:sottomonoidi-generati}:
    in un sottogruppo, oltre ai prodotti di elementi di $t$, ci sono anche i loro inversi ($x^{-1}$, cioè esponente $-1$).
\end{dimostrazione}
\end{concetto}

% ──────────────────────────────────────────────────────────────
\section{Elementi cancellabili}\label{sec:cancellabili}

\dfn[cancellabile]{Elemento cancellabile}{
    Sia $(s, *)$ un semigruppo. Un elemento $a \in s$ si dice
    \begin{itemize}
        \item \textbf{cancellabile a sinistra} se $\forall b, c \in s\,(a * b = a * c \rightarrow b = c)$;
        \item \textbf{cancellabile a destra} se $\forall b, c \in s\,(b * a = c * a \rightarrow b = c)$;
        \item \textbf{cancellabile} se lo è sia a sinistra che a destra.
    \end{itemize}
}

\dfn[traslazioni]{Traslazione sinistra e destra}{
    Sia $(s, *)$ un semigruppo e $a \in s$.
    \begin{enumerate}
        \item \textbf{Traslazione sinistra}: $\sigma_a \colon x \in s \mapsto a * x \in s$.
        \item \textbf{Traslazione destra}: $\delta_a \colon x \in s \mapsto x * a \in s$.
    \end{enumerate}
}

\begin{concetto}
\begin{Proposizione}{Cancellabilità e traslazioni}{traslazioni}
    \begin{prerequisiti}
        \dip{Definizione}{def:cancellabile}
        \dip{Definizione}{def:traslazioni}
        \dip{Proposizione}{prop:forme-iniettivita}
        \dip{Definizione}{def:invertibile}
    \end{prerequisiti}
    Sia $(s, *)$ un semigruppo e $a \in s$.
    \begin{enumerate}
        \item $a$ è cancellabile a sinistra se e solo se $\sigma_a$ è iniettiva.
        \item $a$ è cancellabile a destra se e solo se $\delta_a$ è iniettiva.
        \item Se $s$ è un monoide e $a$ è invertibile, $\sigma_{a^{-1}}$ è l'inversa di $\sigma_a$ e $\delta_{a^{-1}}$ è l'inversa di $\delta_a$.
    \end{enumerate}
\end{Proposizione}
\begin{dimostrazione}
    \begin{caso}{Punto 1}
        $\sigma_a$ iniettiva vuol dire $\forall b, c\,(\sigma_a(b) = \sigma_a(c) \rightarrow b = c)$, cioè $\forall b, c\,(a * b = a * c \rightarrow b = c)$,
        che è la definizione di cancellabile a sinistra.
    \end{caso}
    \begin{caso}{Punto 2}
        Identico, con $\delta_a(b) = b * a$.
    \end{caso}
    \begin{caso}{Punto 3}
        $\sigma_{a^{-1}}(\sigma_a(x)) = a^{-1} * (a * x) = x$ e $\sigma_a(\sigma_{a^{-1}}(x)) = a * (a^{-1} * x) = x$; per $\delta$ è analogo.
    \end{caso}
\end{dimostrazione}
\end{concetto}

\begin{concetto}
\begin{Teorema}{Gli invertibili sono cancellabili}{invertibile-cancellabile}
    \begin{prerequisiti}
        \dip{Definizione}{def:invertibile}
        \dip{Definizione}{def:cancellabile}
    \end{prerequisiti}
    Sia $(m, *, 1_m)$ un monoide. Se $a \in m$ è invertibile a sinistra, allora $a$ è cancellabile a sinistra (e lo stesso a destra).
    In particolare, ogni elemento invertibile è cancellabile.
\end{Teorema}
\begin{dimostrazione}
    Per ipotesi trovo $b \in m$ con $b * a = 1_m$. Se $a * c = a * d$, allora
    \[ c = 1_m * c = b * a * c = b * a * d = 1_m * d = d. \]
    A destra è identico, usando un inverso destro.
\end{dimostrazione}
\end{concetto}

Il viceversa in generale è falso: in $(\mbZ, \cdot, 1)$ ogni $n \neq 0$ è cancellabile, ma solo $-1$ e $1$ sono invertibili.
Diventa vero se il monoide è \textbf{finito}; per dimostrarlo serve il seguente lemma.

\begin{concetto}
\begin{Lemma}{Iniettiva su un insieme finito}{finito-iniettiva}
    \begin{prerequisiti}
        \dip{Definizione}{def:iniettiva}
        \dip{Definizione}{def:suriettiva}
        \dip{Teorema}{teo:comp-iniettive}
        \dip{Teorema}{teo:comp-suriettive}
        \dip{Definizione}{def:funzioni-notevoli}
    \end{prerequisiti}
    Sia $s$ un insieme finito e $f \colon s \to s$ iniettiva. Allora $f$ è suriettiva.
\end{Lemma}
\begin{dimostrazione}
    Per induzione sul numero $n$ di elementi di $s$: dimostriamo che, per ogni $n$, ogni funzione iniettiva di un insieme con
    $n$ elementi in sé stesso è suriettiva.
    \begin{caso}{Passo base: $n = 0$}
        $s = \varnothing$ e la condizione di suriettività ($\forall y \in \varnothing \dots$) è vera a vuoto.
    \end{caso}
    \begin{caso}{Passo induttivo: da $n$ a $n + 1$}
        Supponiamo l'enunciato vero per gli insiemi con $n$ elementi, e sia $s$ con $n + 1$ elementi e $f \colon s \to s$ iniettiva.
        Fisso $a \in s$.
        \begin{caso}{Caso $f(a) = a$}
            Sia $s' = s \setminus \{a\}$, che ha $n$ elementi. Per $x \in s'$ si ha $f(x) \neq f(a) = a$ (iniettività), quindi
            $f_{|s'}$ va da $s'$ in $s'$ ed è iniettiva. Per ipotesi induttiva è suriettiva su $s'$; insieme a $f(a) = a$,
            ogni elemento di $s$ è un'immagine: $f$ è suriettiva.
        \end{caso}
        \begin{caso}{Caso $f(a) = b \neq a$}
            Sia $\tau \colon s \to s$ la funzione che scambia $a$ e $b$ e lascia fissi gli altri elementi: $\tau \circ \tau = \mathrm{id}_s$,
            quindi $\tau$ è biettiva. La funzione $g = \tau \circ f$ è iniettiva (composta di iniettive) e $g(a) = \tau(b) = a$, quindi per
            il caso precedente è suriettiva. Allora $f = \tau \circ g$ (perché $\tau \circ \tau = \mathrm{id}_s$) è composta di suriettive,
            quindi suriettiva.
        \end{caso}
    \end{caso}
\end{dimostrazione}
\end{concetto}

\begin{concetto}
\begin{Teorema}{Monoide finito: invertibile se e solo se cancellabile}{monoide-finito}
    \begin{prerequisiti}
        \dip{Teorema}{teo:invertibile-cancellabile}
        \dip{Proposizione}{prop:traslazioni}
        \dip{Lemma}{lem:finito-iniettiva}
        \dip{Teorema}{teo:unicita-inverso}
    \end{prerequisiti}
    Sia $(m, *, 1_m)$ un monoide finito. Allora $a \in m$ è invertibile se e solo se $a$ è cancellabile.
\end{Teorema}
\begin{dimostrazione}
    \begin{caso}{$(\rightarrow)$ Se $a$ è invertibile, allora è cancellabile}
        Già dimostrato (Teorema~\ref{teo:invertibile-cancellabile}).
    \end{caso}
    \begin{caso}{$(\leftarrow)$ Se $a$ è cancellabile, allora è invertibile}
        $a$ cancellabile $\Rightarrow$ $\sigma_a$ e $\delta_a$ iniettive (Proposizione~\ref{prop:traslazioni}); ma stiamo parlando di un insieme
        finito, quindi per il Lemma~\ref{lem:finito-iniettiva} $\sigma_a$ e $\delta_a$ sono suriettive. Cioè trovo $b, c \in m$ con
        \[ \sigma_a(b) = a * b = 1_m \qquad\qquad \delta_a(c) = c * a = 1_m. \]
        Quindi $a$ ha un inverso destro e uno sinistro: $a$ è invertibile (e $b = c$ per il Teorema~\ref{teo:unicita-inverso}).
    \end{caso}
\end{dimostrazione}
\end{concetto}

% ──────────────────────────────────────────────────────────────
\section{Tavole di Cayley}\label{sec:cayley}

\dfn[cayley]{Tavola di Cayley}{
    Sia $(s, *)$ una struttura con $s$ finito. La \textbf{tavola di Cayley} si ottiene scrivendo gli elementi di $s$, nello stesso ordine,
    su una riga e su una colonna, e mettendo all'incrocio tra la riga di $a$ e la colonna di $b$ il risultato $a * b$.
}

\begin{esempio}{La tavola di $(t, \cdot)$ con $t = \{-1, 0, 1\}$}
    \begin{center}
    \renewcommand{\arraystretch}{1.3}
    \begin{tabular}{c|ccc}
        $\cdot$ & $-1$ & $0$ & $1$ \\ \hline
        $-1$ & \cellcolor{pink!45}$1$ & $0$ & \cellcolor{azure-gradient-2}$-1$ \\
        $0$  & $0$ & \cellcolor{pink!45}$0$ & \cellcolor{azure-gradient-2}$0$ \\
        $1$  & \cellcolor{green-gradient-3}$-1$ & \cellcolor{green-gradient-3}$0$ & \cellcolor{pink!45}$1$ \\
    \end{tabular}
    \end{center}
    \begin{itemize}
        \item \colorbox{green-gradient-3}{Riga di $1$}: ripete la prima riga, quindi $1$ è \textbf{neutro a sinistra};
              \colorbox{azure-gradient-2}{colonna di $1$}: ripete la prima colonna, quindi $1$ è \textbf{neutro a destra}.
        \item La tavola è simmetrica rispetto alla \colorbox{pink!45}{diagonale}: l'operazione è \textbf{commutativa}.
        \item Nella riga (e nella colonna) di $-1$ compare il neutro $1$: $-1$ è \textbf{invertibile} (è l'inverso di sé stesso).
        \item Nella riga di $0$ non compare $1$ e ci sono elementi ripetuti: $0$ non è né invertibile né cancellabile.
    \end{itemize}
\end{esempio}

\begin{concetto}
\begin{Proposizione}{Come si legge una tavola di Cayley}{lettura-cayley}
    \begin{prerequisiti}
        \dip{Definizione}{def:cayley}
        \dip{Definizione}{def:neutro}
        \dip{Definizione}{def:invertibile}
        \dip{Proposizione}{prop:traslazioni}
        \dip{Definizione}{def:commutativa-associativa}
    \end{prerequisiti}
    Sia $(s, *)$ con $s$ finito e sia $a \in s$. Nella tavola di Cayley:
    \begin{enumerate}
        \item $a$ è neutro a sinistra $\leftrightarrow$ sulla riga di $a$ si presentano gli stessi elementi di $s$, nello stesso ordine;
        \item $a$ è neutro a destra $\leftrightarrow$ sulla colonna di $a$ si presentano gli stessi elementi di $s$, nello stesso ordine;
        \item $a$ è simmetrizzabile a sinistra $\leftrightarrow$ nella sua colonna trovo l'elemento neutro;
        \item $a$ è simmetrizzabile a destra $\leftrightarrow$ nella sua riga trovo l'elemento neutro;
        \item $a$ è cancellabile a sinistra $\leftrightarrow$ sulla riga di $a$ ho tutti elementi distinti;
        \item $a$ è cancellabile a destra $\leftrightarrow$ sulla colonna di $a$ ho tutti elementi distinti;
        \item $*$ è commutativa $\leftrightarrow$ la tavola è simmetrica rispetto alla diagonale.
    \end{enumerate}
    (Nei punti 3 e 4 si suppone che $(s, *)$ sia un monoide.)
\end{Proposizione}
\begin{dimostrazione}
    Ricordiamo che la riga di $a$ contiene i valori $a * x$ (cioè $\sigma_a(x)$) e la colonna di $a$ i valori $x * a$ (cioè $\delta_a(x)$), per $x$ che scorre in $s$
    nell'ordine della tavola.
    \begin{caso}{Punti 1 e 2: neutri}
        La riga di $a$ ripete l'intestazione se e solo se $a * x = x$ per ogni $x$, cioè $a$ è neutro a sinistra. Per la colonna, $x * a = x$.
    \end{caso}
    \begin{caso}{Punti 3 e 4: simmetrizzabili}
        Nella colonna di $a$ c'è $1_s$ se e solo se esiste $x$ con $x * a = 1_s$, cioè $a$ ha un inverso sinistro. Nella riga, $a * x = 1_s$: inverso destro.
    \end{caso}
    \begin{caso}{Punti 5 e 6: cancellabili}
        La riga di $a$ ha tutti elementi distinti se e solo se $\sigma_a$ è iniettiva, cioè (Proposizione~\ref{prop:traslazioni}) $a$ è cancellabile
        a sinistra. Per la colonna si usa $\delta_a$.
    \end{caso}
    \begin{caso}{Punto 7: commutatività}
        L'elemento in posizione (riga $x$, colonna $y$) è $x * y$, quello in posizione simmetrica (riga $y$, colonna $x$) è $y * x$:
        coincidono tutti se e solo se $*$ è commutativa.
    \end{caso}
\end{dimostrazione}
\end{concetto}

% ──────────────────────────────────────────────────────────────
\section{Omomorfismi}\label{sec:omomorfismi}

\dfn[omomorfismo]{Omomorfismo tra strutture algebriche}{
    Siano $(s, *)$ e $(s', *')$ strutture algebriche con un'operazione binaria interna. Una funzione $\varphi \colon s \to s'$ si dice
    \textbf{omomorfismo} da $s$ a $s'$ se
    \[ \forall a, b \in s\,\big(\varphi(a * b) = \varphi(a) *' \varphi(b)\big). \]
    Se un omomorfismo è
    \begin{enumerate}
        \item iniettivo, si dice \textbf{monomorfismo};
        \item suriettivo, si dice \textbf{epimorfismo};
        \item biettivo, si dice \textbf{isomorfismo}.
    \end{enumerate}
}

\begin{esempio}{Omomorfismi}
    \begin{itemize}
        \item $m \in \mbZ \mapsto 2m \in \mbZ$ è un monomorfismo di $(\mbZ, +)$ in sé stesso: $2(m + n) = 2m + 2n$.
        \item $m \in \mbZ \mapsto -m \in \mbZ$ è un isomorfismo di $(\mbZ, +)$ in sé stesso.
        \item $\exp \colon x \in \mbR \mapsto e^x \in \mbR^+$ è un isomorfismo da $(\mbR, +)$ a $(\mbR^+, \cdot)$: $e^{x + y} = e^x e^y$.
        \item Se $a$ è un insieme, $x \in \mcP(a) \mapsto a \setminus x \in \mcP(a)$ è un isomorfismo tra $(\mcP(a), \cup)$ e $(\mcP(a), \cap)$:
              è biettiva e, per De Morgan (Proposizione~\ref{prop:demorgan-ins}), $a \setminus (x \cup y) = (a \setminus x) \cap (a \setminus y)$.
    \end{itemize}
\end{esempio}

\begin{concetto}
\begin{Teorema}{L'inversa di un isomorfismo è un isomorfismo}{inverso-isomorfismo}
    \begin{prerequisiti}
        \dip{Definizione}{def:omomorfismo}
        \dip{Teorema}{teo:biettiva-inversa}
    \end{prerequisiti}
    Sia $\varphi \colon s \to s'$ un isomorfismo tra $(s, *)$ e $(s', *')$. Allora $\varphi^{-1}$ è un isomorfismo tra $(s', *')$ e $(s, *)$.
\end{Teorema}
\begin{dimostrazione}
    $\varphi^{-1}$ è biettiva (ha come inversa $\varphi$). Siano $a, b \in s'$:
    \[ \varphi\big(\varphi^{-1}(a) * \varphi^{-1}(b)\big) = \varphi(\varphi^{-1}(a)) *' \varphi(\varphi^{-1}(b)) = a *' b; \]
    applicando $\varphi^{-1}$ a entrambi i membri ho $\varphi^{-1}(a) * \varphi^{-1}(b) = \varphi^{-1}(a *' b)$.
\end{dimostrazione}
\end{concetto}

\begin{concetto}
\begin{Teorema}{Gli epimorfismi conservano le proprietà}{epimorfismi}
    \begin{prerequisiti}
        \dip{Definizione}{def:omomorfismo}
        \dip{Definizione}{def:commutativa-associativa}
        \dip{Definizione}{def:neutro}
        \dip{Definizione}{def:invertibile}
        \dip{Definizione}{def:cancellabile}
    \end{prerequisiti}
    Sia $\varphi \colon s \to s'$ un epimorfismo tra $(s, *)$ e $(s', *')$. Allora in $(s', *')$ vengono conservate associatività,
    commutatività, neutri e simmetrici:
    \begin{enumerate}
        \item se $*$ è commutativa, $*'$ è commutativa;
        \item se $*$ è associativa, $*'$ è associativa;
        \item se $z$ è neutro destro/sinistro di $(s, *)$, allora $\varphi(z)$ è neutro destro/sinistro di $(s', *')$;
        \item se $x$ è simmetrizzabile a destra/sinistra in $(s, *)$, allora $\varphi(x)$ lo è in $(s', *')$;
        \item se $\varphi$ è anche iniettiva (isomorfismo) e $x$ è cancellabile a destra/sinistra, allora $\varphi(x)$ lo è.
    \end{enumerate}
\end{Teorema}
\begin{dimostrazione}
    In tutti i punti, dati elementi di $s'$, li scriviamo come immagini di elementi di $s$ (possiamo perché $\varphi$ è suriettiva).
    \begin{caso}{Punto 1: commutatività}
        Siano $a, b \in s'$ e $x, y \in s$ con $\varphi(x) = a$, $\varphi(y) = b$. Allora
        $a *' b = \varphi(x) *' \varphi(y) = \varphi(x * y) = \varphi(y * x) = \varphi(y) *' \varphi(x) = b *' a$.
    \end{caso}
    \begin{caso}{Punto 2: associatività}
        Siano $a, b, c \in s'$ immagini di $x, y, z$. Allora
        $a *' (b *' c) = \varphi(x) *' \varphi(y * z) = \varphi(x * (y * z)) = \varphi((x * y) * z) = \varphi(x * y) *' \varphi(z) = (a *' b) *' c$.
    \end{caso}
    \begin{caso}{Punto 3: neutri}
        Sia $z$ neutro destro, $a \in s'$ e $x \in s$ con $\varphi(x) = a$. Allora $a *' \varphi(z) = \varphi(x) *' \varphi(z) = \varphi(x * z) = \varphi(x) = a$.
        A sinistra è identico.
    \end{caso}
    \begin{caso}{Punto 4: simmetrici}
        Sia $u$ il neutro di $s$ e $y$ con $x * y = u$. Allora $\varphi(x) *' \varphi(y) = \varphi(x * y) = \varphi(u)$, che per il punto 3 è il neutro
        di $s'$: $\varphi(x)$ è simmetrizzabile a destra. A sinistra è identico.
    \end{caso}
    \begin{caso}{Punto 5: cancellabili}
        Sia $x$ cancellabile a sinistra, $b, c \in s'$ con $\varphi(x) *' b = \varphi(x) *' c$, e $y, z$ con $\varphi(y) = b$, $\varphi(z) = c$. Allora
        $\varphi(x * y) = \varphi(x) *' b = \varphi(x) *' c = \varphi(x * z)$; poiché $\varphi$ è iniettiva, $x * y = x * z$, e poiché $x$ è cancellabile
        a sinistra, $y = z$, quindi $b = c$.
    \end{caso}
\end{dimostrazione}
\end{concetto}
